\naslov{Posplošitve problema lastnih vrednosti}

\podnaslov{Posplošen problem lastnih vrednosti}

Imamo matriki $A, B \in \C^{n \times n}$.
Množico vseh matrik $A - \lambda B$ za $\lambda \in \C$ imenujemo
\pojem{matrični šop}.
Včasih ga označimo z $(A,B)$.
Definiramo lahko karakteristični polinom $p(\lambda) = \det (A - \lambda B)$.
Če je $p = 0$, pravimo, da je šop \pojem{singularen}, sicer pa je
\pojem{regularen}.

Naj bo $(A,B)$ regularen šop.
Če je $A x = \lambda B x$ za $x \ne 0$, je $\lambda$ končna lastna vrednost in
$x$ desni lastni vektor.
Podobno za $y^H A = \lambda y^H B$.
Če pa je $B x = 0$ za nek $x \ne 0$, je $\infty$ lastna vrednost in $x$ desni
lastni vektor; podobno za $y^H B = 0$.
Za regularen šop ima problem $n$ lastnih vrednosti, tj.~ničle karakterističnega
polinoma, ki jim dodamo še toliko neskončnih lastnih vrednosti, da pridemo do
$n$.

\vprasanje{Opiši posplošen problem lastnih vrednosti.}

\begin{izrek}
  Naj bo $(A,B)$ regularen šop.
  \begin{itemize}
  \item Če je $B$ obrnljiva matrika, so lastne vrednosti tega šopa enake lastnim
	vrednostim $B^{-1} A$ oziroma $A B^{-1}$.
  \item $\infty$ je lastna vrednost šopa natanko tedaj, ko je $\det B = 0$.
  \item Če je $A$ nesingularna, potem so lastne vrednosti $(A,B)$ recipročne
	lastnim vrednostim matrike $A^{-1} B$ oz.~$B A^{-1}$.
	Tu je $\infty = \inv{0}$.
  \end{itemize}
\end{izrek}

\vprasanje{Kakšne so lastne vrednosti regularnega šopa $(A,B)$, če je $A$
  obrnljiva?}

\begin{definicija}
  Če sta $U, V$ nesingularni matriki, je šop $(A,B)$ \pojem{ekvivalenten} šopu
  $(UAV, UBV)$.
\end{definicija}

\begin{izrek}
  Ekvivalentna šopa imata enake lastne vrednosti.
  Če je $x$ desni lastni vektor za $(A,B)$, je $V^{-1}x$ desni lastni vektor za
  $(UAV, UBV)$.
  Če je $y$ levi lastni vektor za $(A,B)$, je $U^{-H} y$ levi lastni vektor za
  $(UAV, UBV)$.
\end{izrek}

\vprasanje{Kaj lahko poveš o lastnih vektorjih ekvivalentnih šopov?}

\begin{izrek}
  Za regularen matrični šop $(A,B)$ obstajata taki unitarni matriki $Q,Z$, da je
  $(Q^H A Z, Q^H B Z) = (S,T)$, kjer sta $S$ in $T$ zgornje trikotni matriki.
  Lastne vrednosti $(A,B)$ so potem
  \[
	\lambda_i = \frac{s_{ii}}{t_{ii}}
  \]
  za $t_{ii} \ne 0$ in $\lambda_i = \infty$ za $t_{ii} = 0$.
\end{izrek}

\begin{proof}
  Indukcija na $n$.
  Pri $n = 1$ je trivialno, poglejmo primer $n-1 \leadsto n$.
  Ker je $(A,B)$ regularen šop, ima vsaj eno lastno vrednost $\lambda$ in lastni
  vektor $x$.
  Če je $\lambda$ končna, je $Ax = \lambda B x$, sicer pa je $Bx = 0$.
  V vsakem primeru sta vektorja $Ax$ in $Bx$ kolinearna, torej obstaja $y$ z
  $\norm{y}_2 = 1$, da je $Ax = \alpha y$ in $Bx = \beta y$, ter $(\alpha,
  \beta) \ne (0,0)$.

  Dopolnimo $x$ do unitarne matrike $X = [x X_1]$ in $y$ do unitarne $Y = [y
  Y_1]$.
  Potem je
  \[
	Y^H A X =
	\begin{bmatrix}
	  \alpha & \cdots \\
	  0 & A_1
	\end{bmatrix}
  \]
  in
  \[
	Y^H B X =
	\begin{bmatrix}
	  \beta & \cdots \\
	  0 & B_1
	\end{bmatrix}.
  \]
  Obstajata $Q_1$ in $Z_1$, da je $Q_1^H (A_1 - \lambda B_1) Z_1 = S_1 - \lambda
  T_1$; velja
  \begin{gather*}
	\begin{bmatrix}
	  1 & \\ & Q_1^H
	\end{bmatrix}
	Y^H A X
	\begin{bmatrix}
	  1 & \\ & Z_1
	\end{bmatrix}
	=
	\begin{bmatrix}
	  \alpha & \cdots \\
	  & S_1
	\end{bmatrix}
	= S \\
	\begin{bmatrix}
	  1 & \\ & Q_1^H
	\end{bmatrix}
	Y^H B X
	\begin{bmatrix}
	  1 & \\ & Z_1
	\end{bmatrix}
	=
	\begin{bmatrix}
	  \beta & \cdots \\
	  & T_1
	\end{bmatrix}
	= T. \\
  \end{gather*}
\end{proof}

\vprasanje{Dokaži, da za regularen matrični šop obstaja posplošena Schurova
  forma.}

Za računanje posplošene Schurove forme imamo na voljo QZ algoritem.
Predpriprava je, da z ortogonalno ekvivalenčno transformacijo pretvorimo $(A,B)$
v $(A_1, B_1)$, kjer je $A_1$ zgornje Hessenbergova in $B_1$ zgornje trikotna.
Iščemo torej ortogonalni $Q_1, Z_1$, da velja $A = Q_1 A_1 Z_1$, $B = Q_1 B_1
Z_1$.
V prvem koraku uporabimo algoritem za QR razcep in poiščemo ortogonalno $U$, da
je $U^T B$ zgornje trikotna.
Potem z rotacijami iz leve uničujemo elemente v $A$ (oz.~$U^TA$).
Z vsako rotacijo si pokvarimo en element v $B$, ki ga popravimo z rotacijo z
desne (ta pa nam ne pokvari ustvarjenih ničel v $A$).

\vprasanje{Opiši postopek predpriprave za QZ algoritem.}

Če ima $B$ na diagonali ničlo, jo lahko prestavimo na spodnje desno mesto in
nadaljujemo na podmatriki $(n-1) \times (n-1)$.
Predpostavimo torej, da je $B$ nesingularna.
Matrika $C_0 = A_0 B_0^{-1}$ je zgornja Hessenbergova.
Na njej naredimo korak QR iteracije z dvojnim premikom, da dobimo $C_1 = Q_0^T
C_0 Q_0$, kjer je $Q_0$ matrika iz QR razcepa $N_0 = C_0^2 - (\sigma_1+\sigma_2)
C_0 + \sigma_1 \sigma_2 I$.

QZ iteracija nam iz matrik $A_0, B_0$ da matriki $A_1 = Q_0^T A_0 Z_0$ in $B_1 =
Q_0^T B_0 Z_0$, kjer je $Q_0$ ortogonalna in taka, da je prvi stolpec enak
normiranemu prvemu stolpcu $N_0$, ter $Z_0$ ortogonalna in taka, da je $A_1$
zgornje Hessenbergova in $B_1$ zgornje trikotna.
Po izreku o implicitnem Q je $A_1 B_1^{-1} = Q_0^T A_0 B_0^{-1} Q_0 = C_1$.
Za določitev premika potrebujemo spodnjo desno $2 \times 2$ podmatriko $C_0$,
torej moramo izračunati le inverz spodnje desne $3 \times 3$ podmatrike $B_0$.
Za prvi stolpec $N_0$ potrebujemo tudi inverz vodilne $2 \times 2$ podmatrike
$B_0$.

Prvi stolpec $N_0$ ima tri elemente, torej za redukcijo potrebujemo $3 \times 3$
Householderjevo zrcaljenje.
Ko s tem množimo $A$ in $B$, dobimo grbi; v $A$ je grba velikosti $1 \times 1$,
v $B$ pa velikosti $2 \times 2$.
Potem uničimo grbo v $B$ z dvema zrcaljenjema iz desne, kar nam poveča grbo v
$A$ na velikost $2 \times 2$.
To grbo spravimo nazaj v $1 \times 1$ z zrcaljenjem iz leve, kar nam zopet
prinese grbo v $B$ velikosti $2 \times 2$.
Ko grbo premaknemo do konca, celoten postopek ponavljamo.
QZ iteracija je ekvivalentna inverzni iteraciji na matriki $A B^{-1}$.
Je bolj numerično stabilna kot množenje z inverzom in računanje lastnih
vrednosti, a počasnejša.

\vprasanje{Opiši QZ algoritem. Čemu je ekvivalenten?}

\podnaslov{Kvadratni problem lastnih vrednosti}

Naj bo $Q(\lambda) = \lambda^2 M + \lambda C + K$, kjer so $M$, $C$ in $K$
matrike dimenzij $n \times n$.
Potem je $\det Q(\lambda)$ polinom stopnje manjše ali enake $2n$.
Če ta polinom ni konstantno enak $0$, pravimo, da je problem \pojem{regularen}.

Če je za regularen problem $Q(\lambda) x = 0$, je $\lambda$ lastna vrednost in
$x$ desni lastni vektor.
Podobno je za $y^H Q(\lambda) = 0$ $y$ levi lastni vektor.
V primeru $Mx = 0$ za neničeln $x$ pa pravimo, da je $\infty$ lastna vrednost za
$x$.
Lastne vrednosti so ničle karakterističnega polinoma, ki jih dopolnimo do $2n$ z
neskončnimi.

Klasičen način reševanja je linearizacija.
Če je $\det M \ne 0$, imamo $2n$ končnih lastnih vrednosti, in poiščemo lastne
vrednosti matrike
\[
  \begin{bmatrix}
	0 & I \\
	-M^{-1} K & -M^{-1} C
  \end{bmatrix}.
\]
V primeru $\det M = 0$, pa poiščemo rešitve $(A-\lambda B) u = 0$ za
\begin{align*}
  A =
  \begin{bmatrix}
	C & K \\ K & 0
  \end{bmatrix},
  && B =
	 \begin{bmatrix}
	   -M & 0 \\ 0 & K
	 \end{bmatrix}.
\end{align*}

\vprasanje{Opiši kvadratni problem lastnih vrednosti. Kako ga rešiš?}

% LocalWords:  QZ linearizacija
